\section{Билет 26. Принцип суперпозиции волн. Стоячие волны. Уравнение стоячей волны. Амплитуда стоячей волны. Узлы и пучности.}

\begin{definition}
Важным частным случаем интерференции является наложение двух плоских монохроматических волн с одинаковой амплитудой $A$ и частотой $\omega$, распространяющихся навстречу друг другу. Возникающий колебательный процесс называется \textbf{стоячей волной}. Практически стоячие волны образуются при отражении бегущей волны от границы раздела сред.
Уравнения встречных волн имеют вид:
\begin{equation}
    \xi_1 = A \cos(\omega t - kx), \quad \xi_2 = A \cos(\omega t + kx).
\end{equation}
\end{definition}

\begin{theorem}[Уравнение стоячей волны]
Согласно принципу суперпозиции, результирующее смещение равно:
\begin{equation}
    \xi = \xi_1 + \xi_2 = A \cos(\omega t - kx) + A \cos(\omega t + kx).
\end{equation}
Используя тригонометрическую формулу суммы косинусов $\cos\alpha + \cos\beta = 2\cos\frac{\alpha+\beta}{2}\cos\frac{\alpha-\beta}{2}$, получаем:
\begin{equation}
    \xi(x,t) = 2A \cos(kx) \cos(\omega t). \label{eq:standing_wave}
\end{equation}
Вводя волновое число $k = 2\pi/\lambda$, запишем каноническое \textbf{уравнение стоячей волны}:
\begin{equation}
    \xi(x,t) = \left[ 2A \cos\left(\frac{2\pi x}{\lambda}\right) \right] \cos(\omega t).
\end{equation}
\end{theorem}

\begin{definition}
В уравнении \eqref{eq:standing_wave} множитель в квадратных скобках зависит только от координаты $x$ и представляет собой \textbf{амплитуду стоячей волны} в данной точке:
\begin{equation}
    A_{\text{ст}}(x) = 2A \left| \cos\left(\frac{2\pi x}{\lambda}\right) \right|. \label{eq:standing_amp}
\end{equation}
В отличие от бегущей волны, в стоячей волне амплитуда не переносится в пространстве, а образует стационарную картину чередования максимумов и нулей.
\end{definition}

\begin{property}[Узлы и пучности]
Точки среды, в которых амплитуда колебаний максимальна ($A_{\text{ст}} = 2A$), называются \textbf{пучностями}. Из условия $|\cos(2\pi x/\lambda)| = 1$ находим их координаты:
\begin{equation}
    x_{\text{пучн}} = \pm n \frac{\lambda}{2}, \quad n = 0, 1, 2, \dots \label{eq:antinodes}
\end{equation}
Точки, в которых амплитуда равна нулю ($A_{\text{ст}} = 0$), называются \textbf{узлами}. Из условия $\cos(2\pi x/\lambda) = 0$ получаем:
\begin{equation}
    x_{\text{узел}} = \pm (2n+1) \frac{\lambda}{4}, \quad n = 0, 1, 2, \dots \label{eq:nodes}
\end{equation}
\end{property}

\begin{remark}
\textbf{Свойства стоячей волны:}
\begin{enumerate}
    \item Расстояние между соседними пучностями, а также между соседними узлами равно половине длины волны: $\Delta x = \lambda/2$.
    \item Между двумя соседними узлами все точки среды колеблются \textbf{синфазно}. При переходе через узел фаза колебаний скачкообразно изменяется на $\pi$, то есть соседние сегменты между узлами колеблются в \textbf{противофазе}.
    \item Энергия в стоячей волне не переносится вдоль среды (средний за период поток энергии равен нулю), а происходит её периодическое перераспределение между кинетической и потенциальной формами внутри каждого сегмента.
\end{enumerate}
\end{remark}
